\documentclass[12pt]{article}
\usepackage[T1]{fontenc}
\usepackage[utf8]{inputenc}
\usepackage{lmodern}
\usepackage{textcomp}
\usepackage{enumitem}
\usepackage{geometry}
\geometry{margin=1in}
\begin{document}
\begin{center}
Answer Key - Geography - Graduate Level Assessment 14 \
\small (Answers are ordered exactly as in the question paper.)
\end{center}

\section*{Multiple Choice}
\begin{enumerate}[label=\arabic*.]
\item \textbf{Answer:} C
\\ \textit{Options:} A) China; B) Japan; C) Chile; D) India
\\ \textit{Note:} Chile invested the largest fraction of its GDP in renewable energy in 2015 due to its robust policies promoting solar and wind energy, driven by the country's abundant natural resources and a commitment to sustainable development.
\item \textbf{Answer:} D
\\ \textit{Options:} A) 50\%; B) 60\%; C) 70\%; D) 80\%
\\ \textit{Note:} The correct answer of 80\% reflects the strong value placed on free media as essential to democracy and civil liberties among Americans, highlighting a widespread belief in the importance of media independence from government influence.
\item \textbf{Answer:} D
\\ \textit{Options:} A) shortly after World War II; B) in the 1960 s; C) around the time of World War I; D) in the early nineteenth century
\\ \textit{Note:} Most Latin American countries achieved independence in the early nineteenth century due to a combination of Enlightenment ideals, the influence of the American and French Revolutions, and the weakening of Spanish and Portuguese colonial powers resulting from Napoleonic Wars.
\item \textbf{Answer:} C
\\ \textit{Options:} A) Brazil; B) China; C) Mexico; D) Iran
\\ \textit{Note:} As of 2019, Mexico had the lowest life expectancy among the options due to a combination of factors including socioeconomic disparities, healthcare access issues, and high rates of violence that impact overall health outcomes.
\item \textbf{Answer:} C
\\ \textit{Options:} A) Diarrheal diseases; B) Diabetes; C) Dementia; D) Road injuries
\\ \textit{Note:} Dementia causes more deaths globally each year due to its prevalence as a leading cause of mortality, particularly among the aging population, highlighting the significant health challenges associated with neurodegenerative diseases.
\end{enumerate}

\section*{True / False}
\begin{enumerate}[label=\arabic*.]
\setcounter{enumi}{5}
\item \textbf{Answer:} True
\\ \textit{Options:} A) True; B) False
\\ \textit{Note:} According to the passage: bilingual primary and secondary education
\item \textbf{Answer:} True
\\ \textit{Options:} A) True; B) False
\\ \textit{Note:} According to the passage: the Press Complaints Commission
\item \textbf{Answer:} True
\\ \textit{Options:} A) True; B) False
\\ \textit{Note:} According to the passage: Coarse sediments
\item \textbf{Answer:} False
\\ \textit{Options:} A) True; B) False
\\ \textit{Note:} The correct answer is: force newspapers to print government rebuttals to stories to which the provincial cabinet objected
\item \textbf{Answer:} True
\\ \textit{Options:} A) True; B) False
\\ \textit{Note:} According to the passage: Christianity
\end{enumerate}

\section*{Long Form Response}
\begin{enumerate}[label=\arabic*.]
\setcounter{enumi}{10}
\item \textbf{Answer:} Santa Clara, California.
\\ \textit{Note:} Expected answer: Santa Clara, California.
\item \textbf{Answer:} to identify and exact settlements
\\ \textit{Note:} Expected answer: to identify and exact settlements
\end{enumerate}

\vfill
\noindent\textit{End of Answer Key}
\end{document}